\section{Related Literature}\label{sec:literature}

\subsection{Early Empirical Evidence}

A small but growing body of work investigates employment shifts associated with generative AI exposure. We review the most closely related studies to ours before turning to null findings and to the measurement of AI exposure.

\cite{brynjolfsson2025} provide large-scale evidence from the United States. Using private payroll records from ADP covering 3.5 to 5 million workers per month between January 2021 and September 2025, they estimate a Poisson regression event study with firm-by-quintile and firm-by-time fixed effects. Workers aged 22--25 in the most AI-exposed occupational quintile experienced a decline of approximately 15 log points in employment relative to the least-exposed quintile. The effect is concentrated in occupations exposed to automation, as classified by the \cite{handa2025} decomposition. Occupations exposed primarily to augmentation show no consistent pattern. Wages exhibit little divergence over the same period, consistent with downward nominal rigidity.

\cite{lodefalk2026} replicate this design in Sweden using employer-employee data covering January 2019 through June 2025. Their difference-in-differences specification includes employer-by-quartile and employer-by-month fixed effects, with AI exposure measured by the DAIOE generative AI index of \cite{engberg2024} at the four-digit SSYK level. Workers aged 22--25 in the most-exposed quartile show a 5.5 percent employment decline by the first half of 2025. The age gradient is monotonic: the youngest workers decline most, while the oldest (50 and above) gain 1.3 percent. A separate analysis of job postings from Platsbanken attributes part of the decline in vacancies to monetary tightening following the Riksbank rate hike of April 2022, rather than to AI. Pre-trend joint tests reject the null of parallel trends, but Lodefalk et al.\ argue that the pre-trend slope for the 22--25 group is essentially noise and that any linear extrapolation would bias the estimate toward zero rather than away from it.

\cite{teeselink2025} uses Revelio Labs data from the United Kingdom and finds that a one standard deviation increase in LLM exposure reduces employment by 0.3 percent. The effects are concentrated in junior positions and in high-wage occupations. \cite{teeselink2026} analyze job posting data from 39 countries and find that AI exposure reduces job postings on average. Countries with stricter employment protection exhibit larger hiring reductions, consistent with firms adjusting on the hiring margin when separation costs are high.

Reported magnitudes vary across countries and specifications: Brynjolfsson et al.\ report a 16 log-point relative decline in the US, Lodefalk et al.\ a 5.5 percent decline in Sweden, and Teeselink a 0.3 percent decline per standard deviation in the UK. These figures use different units, baselines, and specifications and should not be read as directly comparable magnitudes.


\subsection{Null and Mixed Results}

Not all studies find displacement. \cite{humlum2026} study Denmark using population-wide register data combined with two survey rounds of approximately 25,000 responses each. Their difference-in-differences design compares workers who adopted AI chatbots to non-adopters. They find a precise null on earnings and no employment effects; at the workplace level, total hours at encouraged-use workplaces are precise enough to rule out effects larger than 2 percent. However, adopters who also switched occupations on average moved to occupations with higher AI exposure and experienced faster earnings growth. This represents an increase in full-time-equivalent occupational mobility, transitioning into higher-wage, higher-AI-exposure occupations.

In their Appendix D.2, \citet{humlum2026} use the same 22--25 early-career age definition as \citet{brynjolfsson2025} and document aggregate declines in early-career employment in several of their exposed occupations in Denmark, including software development, customer support, legal professions, and marketing. The Danish aggregate pattern thus mirrors the US pattern. Their firm-level difference-in-differences then shows that these declines are not driven by firms adopting AI chatbots: estimates rule out effects larger than one third of a percentage point on the early-career employment share at adopting workplaces. They suggest that changes in the macroeconomic environment or anticipatory hiring freezes unrelated to actual adoption may be responsible.

Finnish evidence reinforces the null findings. \cite{kauhanen2026} use population-level register data and replicate the \cite{brynjolfsson2025} methodology. They find no systematic displacement effects and attribute observed employment trends to demographic shifts rather than AI exposure. Earlier Finnish analyses by \cite{kauhanen2024, kauhanen2025} using a C-AIOE exposure measure likewise found null effects.

Measurement, institutional context, and adoption timing differ across these studies. \cite{humlum2026} identify adopters through surveys; \cite{brynjolfsson2025} and \cite{lodefalk2026} rely on occupation-level exposure indices. If AI displaces through reduced firm entry and exit, individual-level adoption designs may miss the channel. Nordic labor markets share strong employment protection and centralized bargaining that may buffer or delay displacement.


\subsection{AI Exposure Measurement}

Empirical studies of AI and employment use occupation-level measures that map AI capabilities to task content. The task framework underlying much of this work is set out in \citep{autor2003skill, acemoglu2011, acemoglu2018race, acemoglu2019automation}. The measures in current use differ in whether they capture potential task exposure, revealed usage, or historical AI progress on the abilities a job requires.

\cite{eloundou2024} develop the most widely used framework for LLM exposure. Building on the O*NET task database, they employ human annotators and GPT-4 to rate each task on whether access to an LLM (or LLM-powered software) would reduce completion time by at least 50 percent. The resulting measures distinguish between whether this reduction of at least 50 percent would occur with direct LLM exposure or with complementary software. ($E_1$) measures the share of an occupation's that are directly exposed, while ($E_2$) measures the share of tasks exposed through LLM and complementary software. Higher-wage occupations score higher on exposure, reversing the pattern observed with earlier waves of automation. The GPT-4 ``beta'' score, defined as $E_1 + 0.5 \times E_2$, has become a standard measure in subsequent empirical work.

\cite{handa2025}, or the Anthropic Economic Index, measures exposure through data on actual usage. They use a privacy-preserving classifier to map approximately four million Claude.ai conversations to O*NET tasks, dividing each task's conversation count by the number of occupations it is associated with before aggregating to the occupation level. They decompose usage into automation (directive task delegation and feedback-loop conversations with minimal human involvement) and augmentation (collaborative task iteration, learning, and validation), reporting 43 percent automation and 57 percent augmentation across all conversations. This decomposition is important because the two modes may have opposite labor market implications: automation substitutes for worker effort, while augmentation complements it. AI use is highly concentrated by occupation: only about 4 percent of occupations show usage in at least three-quarters of their tasks and about 36 percent in at least a quarter, with Computer and Mathematical occupations alone accounting for 37 percent of all observed conversations.

\cite{felten2021} take a different approach, linking AI application benchmarks across ten categories (including language modeling, image recognition, and speech recognition) to the occupational abilities listed in O*NET. Their AI Occupational Exposure (AIOE) score is a weighted sum of AI progress on the abilities required by each occupation. This measure has been extended to a dynamic version (DAIOE) that tracks AI capability growth over time. \cite{engberg2024} construct a DAIOE variant emphasizing generative AI applications, which is used by \cite{lodefalk2026} in the Swedish analysis described above.

A deeper conceptual concern is that standard exposure indices aggregate task-level scores using linear formulas, implicitly treating tasks as separable. \cite{gans2025} argue that in many occupations, tasks are quality complements rather than independent: output is multiplicative across tasks, as in an O-ring production function \citep{kremer1993}. Under task complementarity, automating some tasks frees the worker to focus on the remaining ones, raising their quality and potentially increasing labor income. The relevant object is then not average task exposure but the structure of bottleneck tasks and how automation reshapes worker time around them. \cite{imas2026} extend this logic, arguing that job dimensionality (the number of distinct tasks in a job) determines whether AI augments or displaces. Low-dimensionality jobs built around a small number of core tasks face higher displacement risk even if their average exposure score is low, while high-dimensionality jobs with many complementary tasks are more likely to be augmented. Exposure indices that average across tasks will overstate displacement risk for complex jobs and understate it for narrow ones.

For Norway specifically, existing evidence documents growing AI adoption across the economy. \cite{kompetansebehovsutvalget2026} report that AI-related job postings nearly doubled between 2023 and 2024, concentrated in IT but spreading to other sectors.
